\chapter{Many-electron Wave Functions and Operators}\label{c2}
\section{Spinors}\label{c2s1}
A complete description of an electronic state needs its spatial as well as spin
wavefunction. Being a spin $1/2$ particle, an electron's spin state can be 
described by one of
\begin{equation}\label{c2e1}
\ket{+} = \begin{bmatrix} 1 \\ 0 \end{bmatrix}; \;
\ket{-} = \begin{bmatrix} 0 \\ 1 \end{bmatrix}
\end{equation}
If $\psi_+(\vec{r})$ is the spatial wave function in the spin up state and 
$\psi_-(\vec{r})$ in the spin down state then the electron's state can be
expressed by 
\begin{equation}\label{c2e2}
\ket{\chi} = \psi_+(\vec{r})\ket{+} + \psi_-(\vec{r})\ket{-} =
\begin{bmatrix}\psi_+(\vec{r}) \\ \psi_-(\vec{r}) \end{bmatrix}.
\end{equation}
The probability of finding an electron in a volume element $d\vec{r}$ about the
point $\vec{r}$ is 
\[
\rho(\vec{r})d\vec{r} = (|\psi_+(\vec{r})|^2 + |\psi_-(\vec{r})|^2)d\vec{r}.
\]
When the spin and spatial degrees of freedom are decoupled we can write the 
spinor as
\begin{equation}\label{c2e3}
\ket{\chi} = \psi(\vec{r})\begin{bmatrix}c_1 \\ c_2\end{bmatrix}
\end{equation}
where $c_1, c_2 \in \mathbb{C}$. This is the most common situation in quantum
chemistry. 

Here are the expressions for common operators acting on spinors.
\begin{eqnarray}
p\ket{\chi} &=& -i\hbar\nabla\psi_+\ket{+} - i\hbar\nabla\psi_-\ket{-} = 
-i\hbar\begin{bmatrix}\nabla\psi_+ \\ \nabla\psi_-\end{bmatrix}
\label{c2e4} \\
L\ket{\chi} &=& (L\psi_+)\ket{+} + (L\psi_-)\ket{-} = 
\begin{bmatrix}L \psi_+ \\ L \psi_-\end{bmatrix}
\label{c2e5} \\
S\ket{\chi} &=& \psi_+S\ket{+} + \psi_-S\ket{-}
\label{c2e6}
\end{eqnarray}

In case we have an operator like $H_{SO}=\vec{L}\cdot\vec{S}$,
then we use the Pauli matrix representation for 
\begin{equation}\label{c2e7}
\vec{S} = \frac{\hbar}{2}\left(\sigma_x\uv{x} + \sigma_y\uv{y} +
\sigma_z\uv{z}\right)
\end{equation}
so that
\begin{equation}\label{c2e8}
H_{SO}=\frac{\hbar}{2}\left(L_x\sigma_x + L_y\sigma_y + L_z\sigma_z\right).
\end{equation}
Since
\begin{eqnarray}
\sigma_x &=& \begin{bmatrix}0 & 1 \\ 1 & 0\end{bmatrix} \label{c2e9} \\
\sigma_y &=& \begin{bmatrix}0 & i \\ -i & 0 \end{bmatrix} \label{c2e10} \\
\sigma_z &=& \begin{bmatrix}1 & 0 \\ 0 & -1 \end{bmatrix}\label{c2e11} 
\end{eqnarray}
so that
\begin{equation}\label{c2e12}
H_{SO} = \frac{\hbar}{2}\begin{bmatrix}
L_z & L_x + iL_y \\
L_x - iL_y & -L_z
\end{bmatrix}
\end{equation}
and
\begin{equation}\label{c2e13}
H_{SO}\ket{\chi} = \frac{\hbar}{2}\begin{bmatrix}
L_z & L_x + iL_y \\
L_x - iL_y & -L_z
\end{bmatrix}\begin{bmatrix}\psi_+(\vec{r}) \\ \psi_-(\vec{r})\end{bmatrix}.
\end{equation}

\section{The molecular hamiltonian}
Consider a molecule with $N$ electrons at positions $\vec{r}_1, \ldots, 
\vec{r}_N$, $M$ nuclei with masses $M_1, \ldots, M_M$, atomic numbers $Z_1, 
\ldots, Z_M$ and at positions $\vec{R}_1, \ldots, \vec{R}_M$. The complete,
non-relativistic hamiltonian for the molecule, in atomic units, is
\begin{equation}\label{c2e14}
H = -\frac{1}{2}\sum_{i=1}^N\nabla_i^2 - 
\sum_{i=1}^N\sum_{a=1}^M\frac{Z_a}{R_{ia}} +
\sum_{i=1}^N\sum_{j=i+1}^N\frac{1}{r_{ij}} -
\sum_{a=1}^M\frac{1}{2M_a}\nabla_a^2 +
\sum_{a=1}^M\sum_{b=a+1}^M\frac{Z_aZ_b}{R_{ab}},
\end{equation}
where $r_{ij} = |\vec{r}_i - \vec{r}_j|, r_{ia} = |\vec{r}_i - \vec{R}_a|$ and
$R_{ab} = |\vec{R}_a - \vec{R}_b|$.

Since the nuclei are much more massive than the electrons, substantial 
simplification is achieved by treating them as stationary while studying the 
electronic properties. This is called the \emph{Born-Oppenheimer} approximation
under which the molecular hamiltonian is 
\begin{equation}\label{c2e15}
H_{el} = -\frac{1}{2}\sum_{i=1}^N\nabla_i^2 - 
\sum_{i=1}^N\sum_{a=1}^M\frac{Z_a}{R_{ia}} +
\sum_{i=1}^N\sum_{j=i+1}^N\frac{1}{r_{ij}} -
\end{equation}
Its eigenvalue equation is
\begin{equation}\label{c2e16}
H_{el}\Phi(\vec{r}; \vec{R}) = \mathscr{E}(\vec{R})\Phi(\vec{r}; \vec{R}),
\end{equation}
where $\vec{r} = (\vec{r}_1, \ldots, \vec{r}_N)$ and $\vec{R} = (\vec{R}_1,
\ldots, \vec{R}_M)$. The eigenvalue and the eigenfunction have $\vec{R}$ as a
parameter. We are essentially assuming that the nuclei are at rest and the 
system moves on the potential energy surface
\begin{equation}\label{c2e17}
\mathscr{E}_{nucl} = \sum_{a=1}^M\sum_{b=a+1}^M\frac{Z_aZ_b}{R_{ab}}.
\end{equation}

